\chapter{Every rung}
\label{ch:rungs}

The hypothesis is stated for every rung. At rung $n$ the clock $E_1$ is replaced by the volume element $\Gamma_n$ in the whole Clifford
algebra of the $3n$ temporal directions,
\[
\sigma_n:\ E_1\mapsto\Gamma_n,\ \ X\mapsto X\ (X\in P),\qquad \gn=\mathrm{span}\{\half XY,\ \half X\Gamma_n\},\qquad G_n=\sigma_n(T_n),
\]
the spatial triple of the four-plane carries the mass and no Yang--Mills connection, and the spatial directions of the other constituents
carry the Yang--Mills algebra $\mathfrak s_n$ of (A4). This chapter determines what commutes with the Lorentz algebra at rung $n$, the
algebra $\gn$ and its commutant, the generator of mass, and the structure of the spinor module $\R^{8^n}$ for $n=1,2,3$. Throughout, the
Lorentz action is the spectator action of \cref{def:hyp}.

\section{What commutes with the Lorentz algebra}

\begin{theorem}[The commutant is a Clifford algebra]
\label{thm:commutant}
\begin{enumerate}[label=(\roman*),leftmargin=2.2em,itemsep=1pt]
\item $\Gamma_n$, the $3n-1$ elements of $P$ and the $3n-3$ elements of $\cN$ are $6n-3$ anticommuting elements that commute with $\lint$;
$3n$ of them have square $+\id$ and $3n-3$ have square $-\id$. They generate a Clifford algebra $\Cl(3n,3n-3)$, of dimension $2^{6n-3}$.
\item This algebra is the whole commutant of $\lint$ in the matrix algebra of $\R^{8^n}$.
\item $\Omega=\pm\Gamma_n\prod_P\prod_{\cN}$ is its volume element and is central in it.
\end{enumerate}
\checks{cf-na, cf-nb, gm-1, cf-2i, for $n=1,2,3$}
\end{theorem}

\begin{proof}
(i) $\Gamma_n$ anticommutes with every generator, and an element of $P$ or $\cN$ is a generator other than $E_1,E_4,E_5,E_6$; so the $6n-3$
elements anticommute pairwise, and each anticommutes with $E_1,E_4,E_5,E_6$ and therefore commutes with their bivectors, which span
$\lint$. The squares are $\Gamma_n^2=\id$, $X^2=\id$ for $X\in P$ and $X^2=-\id$ for $X\in\cN$. With the convention of \cref{sec:conv}, in
which the first index counts generators of square $+\id$, $\Cl(p,q)$ with $p-q=3$ is a full matrix algebra over $\C$: here
$\Cl(3n,3n-3)\cong M_{2^{3n-2}}(\C)$, a simple algebra of real dimension $2^{6n-3}$, and its representation on $\R^{8^n}$ is faithful.
(ii) The algebra generated by $\lint$ is $\Cl^0(1,3)\otimes\id\cong M_2(\C)$, acting on the first factor of $\R^{8^n}=\R^8\otimes\R^{8^{n-1}}$.
By the proof of \cref{thm:clifford1}, $\R^8$ is twice the simple module $\C^2$, so $\R^{8^n}$ is $2\cdot8^{n-1}$ copies of it, and the
commutant is $M_{2\cdot8^{n-1}}(\C)$, of real dimension $2\,(2\cdot8^{n-1})^2=2^{6n-3}$. The algebra of (i) lies in the commutant and has
the same dimension, so it is the whole commutant.
(iii) The product $\Gamma_n\prod_P\prod_\cN$ contains $E_1,E_4,E_5,E_6$ once and every other generator twice, so it is $\pm\Omega$. The volume
element of a Clifford algebra with an odd number of generators commutes with each generator.
\end{proof}

In the field reading the four-plane therefore enters what commutes with its Lorentz algebra through one element, the volume element
$\Gamma_n$, in place of the clock. What remains is a space of signature $(3n,3n-3)$: the $3n$ temporal directions with the clock replaced,
and the spatial directions of the other constituents. The two families of generators anticommute, $XN=-NX$ for $X$ in $(\Gamma_n;P)$ and $N$
in $\cN$, so the commutant is the graded tensor product of their Clifford algebras; the even elements of either factor, in particular the
temporal rotations, commute in the ordinary sense with all of the other.

\section{The transported algebra}

Let $W_0$ be the product of all $6n-4$ internal generators: the surplus pair and all six generators of each further constituent.

\begin{theorem}[The hypothesis at rung $n$]
\label{thm:everyrung}
\begin{enumerate}[label=(\roman*),leftmargin=2.2em,itemsep=1pt]
\item $\gn$ is the rotation algebra of the $3n$ anticommuting symmetric involutions $(\Gamma_n;P)$, so $\gn\cong\so{3n}$. Its
$3n(3n-1)/2$ generators are antisymmetric matrices that commute with $\lint$, so they keep the four kinetic forms.
\item The two readings are conjugate. $\mathcal G_n:=\exp(\tfrac\pi4\Gamma_nE_1)$ is orthogonal, sends the clock $E_1$ to $\Gamma_n$ and fixes
$P$; on $\Cl(V^+_n)$ the map $\sigma_n$ is conjugation by $\mathcal G_n$, and $\mathcal G_n\,\kn\,\mathcal G_n^{-1}=\gn$.
\item The two readings differ on the boosts: the elements $XE_1$ of $\kn$ do not commute with them.
\item Every vector of $\R^{8^n}$ lies in a spinor representation of $\so{3n}$, and the Casimir of $\gn$ is the scalar $-\tfrac{3n(3n-1)}8$.
\item $U:=\exp(-\tfrac\pi4W_0)$ is orthogonal, commutes with $\lint$, $\Omega$ and $\Gamma_n$, fixes every $XY$ and sends $X\Gamma_n$ to
$\pm X\Omega$: $U\,\gn\,U^{-1}=\mathrm{span}\{\half XY,\half X\Omega\}$. At $n=1$ the two algebras are equal.
\item The commutant of $\gn\oplus\lint$ is generated by $\Omega$ and $\cN$: it is $\C_\Omega\otimes\Cl(0,3n-3)$, of dimension $2^{3n-2}$.
\end{enumerate}
\checks{cf-nc, cf-nd, cf-ne, cf-2i, er-na to er-nj, er-nr, for $n=1,2,3$; (vi) at $n=2$, with the inclusion at $n=3$}
\end{theorem}

\begin{proof}
(i) to (iii): $\Gamma_n$ is a symmetric involution that anticommutes with every generator, so $(\Gamma_n;P)$ satisfy the relations of
$\Cl(3n,0)$ and commute with the bivectors of the four-plane. $\Gamma_nE_1$ is antisymmetric with square $-\id$, commutes with $P$ and
anticommutes with $E_1$, so $\exp(\tfrac\pi4\Gamma_nE_1)$ turns $E_1$ into $\Gamma_n$ and fixes $P$.
(iv) $\R^{8^n}$ is a module of the algebra $\Cl(3n,0)$ generated by $(\Gamma_n;P)$, hence a sum of its irreducible modules, on each of which
$\so{3n}$ acts by a spinor representation. Each generator $\half XY$ squares to $-\tfrac14\id$, which gives the Casimir.
(v) $W_0$ is a product of $6n-4$ anticommuting generators, $3n-1$ of square $+\id$ and $3n-3$ of square $-\id$, so $W_0^2=-\id$ and
$W_0^\top=-W_0$; it anticommutes with each of its factors and commutes with $E_1,E_4,E_5,E_6$. Hence $UXU^{-1}=XW_0$ for $X\in P$, and
$\Gamma_n=\pm\Omega W_0$ gives $U(X\Gamma_n)U^{-1}=\pm X\Omega$.
(vi) By \cref{thm:commutant} the commutant of $\lint$ is spanned by the products of a blade of $(\Gamma_n;P)$ and a blade of $\cN$. A
rotation $XY$ commutes with every blade of $\cN$, and a blade of $(\Gamma_n;P)$ commutes with all rotations only if it is $1$ or the volume
element $G_n=\Gamma_n\prod_P$, an identity that holds at every $n$ by the definition of $\sigma_n$. So the commutant of $\gn$ and $\lint$ is
$\mathrm{span}\{1,G_n\}\cdot\Cl(\cN)$, and by \cref{thm:commutant}(iii) $\Omega=\pm G_n\,\omega_\cN$, with $\omega_\cN$ invertible.
\end{proof}

\begin{remark}[Conditional on the reading]
\label{rem:conditional}
The displacement reading has $\kn$, the rotations of the $3n$ temporal directions with the clock among them; it moves the clock, and with it
the four-plane and the boosts. The field reading has $\gn$, which leaves the four-plane and its Lorentz algebra where they are. The two are
one algebra seen through the orthogonal map $\mathcal G_n$, which is not in the commutant of $\lint$ and moves the fixed Lorentz embedding.
On the spinors the exclusion of \PI{Prop.~3.6} is this: the same abstract algebra $\so{3n}$ acts in both readings, and only one of the two
actions is compatible with a fixed Lorentz algebra.
\end{remark}

\begin{remark}[The first form]
\label{rem:firstform}
A first form of the hypothesis replaces $E_1$ by $\Omega$ in the rotations only, with the algebra $\mathrm{span}\{\half XY,\half X\Omega\}$. By
(v) it is conjugate to $\gn$ by an orthogonal map that commutes with $\lint$, $\Omega$ and $\Gamma_n$, and equal to it at one constituent.
Statements that do not depend on such a conjugation hold in both forms; several checks of this chapter, those labelled \chk{er} and
\chk{sr}, are made in the first form.
\end{remark}

By \cref{thm:commutant,thm:everyrung} all that commutes with the Lorentz algebra is built from two pieces: the transported Clifford algebra of
the temporal directions, which holds the Higgs vectors, the Yang--Mills algebra and the generator of mass, and the Clifford algebra of the
spatial directions of the other constituents, which together with $\Omega$ is exactly what commutes with the Yang--Mills and the Lorentz
algebras. By (A4) its antisymmetric part $\mathfrak s_n$ is gauged as well. The $3n$ directions of $\gn$ are $\Gamma_n$ with the surplus pair
and the temporal triples of the other constituents, so $\gn$ contains $\so3\oplus\cdots\oplus\so3$, one summand for each constituent: the
image of $\so3_+$ for the first and the plain temporal rotations for the others.

\section{The generator of mass}
\label{sec:test}

By (A3) the generator of mass at rung $n$ is the image of the temporal volume of all $3n$ directions. The following theorem tests this choice
at rungs one, two and three, together with the image of the clock vector, which is the Higgs vector of the first rung.

\begin{theorem}[The generator of mass at rungs one, two and three]
\label{thm:massrung}
For $n=1,2,3$:
\begin{enumerate}[label=(\roman*),leftmargin=2.2em,itemsep=1pt]
\item $G_n=\sigma_n(T_n)=\Gamma_n\prod_P=-\gamma^0S_n=-\Omega\,\omega_\cN$: the temporal volume of the rung is replaced by minus the clock
times the spatial volume of the rung.
\item $G_n$ commutes with $\gn$ and with $\lint$: all of $\so{3n}$ is kept.
\item $G_n^2=(-1)^{3n(3n-1)/2}\,\id$, which is $-\id$ at rungs one and two and $+\id$ at rung three. At rungs one and two $G_n$ is
antisymmetric and generates a rotation; at rung three it is a symmetric involution and generates none.
\item The image of the clock vector, tested as a vacuum. With the value $\kappa\Gamma_n$ the rotations $XY$ of $P$ are unbroken and the $3n-1$
generators $X\Gamma_n$ are broken: $\so{3n}\to\so{3n-1}$. The kernel $\kappa\Gamma_n\omega_S=\pm\kappa\,T_n\omega_\cN$ is antisymmetric,
Lorentz invariant and invariant under $\so{3n-1}$.
\item The Lorentz-invariant mass kernels that are invariant under all of $\gn$ are precisely the matrices $K=\gamma^0C$ with $C$ in the
commutant of \cref{thm:everyrung}(vi) and $K^\top=-K$. Each blade of the commutant gives a symmetric or an antisymmetric $\gamma^0C$, and
$0$, $8$ and $72$ of them are antisymmetric at $n=1,2,3$; they span the kernel space, of these real dimensions. At rung two all eight commute
with $G_2$; at rung three forty commute with $G_3$ and thirty-two anticommute. Invariance under $\gn\oplus\mathfrak s_n$ as well leaves
spaces of dimension $0$, $0$ and $2$.
\item The image of the temporal volume of the first constituent alone is $\sigma_n(T_1)=\Gamma_nJ=-\Omega\,\Gamma'$, with $\Gamma'$ the
volume element of the other constituents. Its square is $-\id$ at each of the three rungs, and within $\gn$ it commutes only with $\so3\oplus\so{3n-3}$.
\end{enumerate}
\checks{cf-nf, cf-ng, cf-nh, cf-nl, ps-9}
\end{theorem}

The sign in (iii) holds for every $n$, since $\sigma_n$ is a homomorphism and $T_n^2=(-1)^{3n(3n-1)/2}$: the generator of mass is a rotation
for $n\equiv1,2$ and an involution for $n\equiv3,0$ modulo four. This is the period four of \PI{Thm~9.62}(ii), where $T_n$ is an
antisymmetric form for $n\equiv1,2$ and a symmetric one for $n\equiv3,0$.

The test thus has a definite outcome. With the volume of all $3n$ temporal directions the generator of mass keeps all of $\so{3n}$. At rungs
one and two it is a rotation and has a rate; at rung three it is an involution, which grades the module and gives no rest rotation. For a
rotation at rung three three candidates remain, and none is selected here: the image of $T_1$, which is a rotation at every rung examined
and keeps only $\so3\oplus\so{3n-3}$; the pseudoscalar $-\Omega$ alone, which is central in the commutant at every rung and acts on the field
strength as $G$, but for $n\ge2$ is not the image of an element of $\Cl(V^+_n)$; or no rotation, the mass at rung three being carried by
kernels only.

\begin{remark}[The generator of mass and the binding root]
\label{rem:Gn}
By (i), $G_n=-\Omega\,\omega_\cN$, the pseudoscalar of the four-plane times the spatial volume of the other constituents. On an embedded
two-form, $G_nF_X=-\Omega F_X\omega_\cN=F_{GX}\,\omega_\cN$ in the ambient algebra. The right-hand side is the product of a two-form of the
four-plane with the disjoint volume $\omega_\cN$, of grade $3n-1$, and not itself a two-form; since $\omega_\cN$ commutes with $F_{GX}$, the
identity becomes an equality of operators on each eigenspace of $\omega_\cN$. At rung two $\omega_\cN$ is a symmetric involution, its
eigenspaces $\omega_\cN=\pm1$ are real, and on them $G_nF_X=\pm F_{GX}$: the generator of mass acts as $\pm G$, the binding rotation with
one sense on each (\cref{thm:rung2}). At rung three $\omega_\cN^2=-\id$; its eigenspaces exist after complexification, with the eigenvalues
$\pm i$, and on them the induced operator is $\pm iG$, of square $+\id$, in agreement with \cref{thm:massrung}(iii). The factor
$\omega_\cN$ lies in the Clifford algebra of $\cN$, and $\omega_\cN^\top=(-1)^{k(k+1)/2}\omega_\cN$ for its $k=3n-3$ generators of square
$-\id$: it is symmetric, and not in $\mathfrak s_2$, at rung two, and antisymmetric, an element of $\mathfrak s_3$, at rung three.
\checks{cf-3m}
\end{remark}

\begin{remark}[The working model, at a test vacuum]
\label{rem:recovered}
If a nonzero vacuum value of the clock vector is imposed, as in the test (iv), with the vector Higgs kinetic term, the massive components at
rung $n$ are $3n-1$ fields in the vector of the unbroken $\so{3n-1}$, whose connection is the massless part. This is the setting of
\cref{ch:bcoupling}, $N'$ fields with the copy symmetry $\so{N'}$ and its Yang--Mills field, with $N'=3n-1$: the copies are the directions of
$P$. At the first rung $N'=2$, the surplus pair, and the connection is abelian, the case computed in \cref{ch:quantum}. Above the first rung
this test vacuum is not the one that (A3) assigns.
\end{remark}

\section{The first rung}

$\mathfrak g_1=\half\,\mathrm{span}\{J,e_2\Omega,e_3\Omega\}$ is the algebra of Pauli and G\"ursey, the commutant of $\mathfrak g_1\oplus\lint$
is $\C_\Omega$, and $\R^8$ is one Weyl doublet. The generator of mass is $G_1=-\Omega$, and no nonzero mass kernel is invariant
(\cref{thm:clifford1,thm:massgen,prop:twojobs,prop:onedoublet}).
\checks{cf-1a to cf-1h, cf-1l, er-1k, er-1l}

\section{The second rung}
\label{sec:rung2}

Write $a=E_2$ and $b=E_3$ for the surplus pair and $c,d,e=E_7,E_8,E_9$ for the temporal triple of the second constituent, as in
\PI{Def.~8.56}, and $I_5=abcde$. Here $\cN=\{N_1,N_2,N_3\}=\{E_{10},E_{11},E_{12}\}$, and
\[
S_i=\half N_jN_k\quad(i,j,k\ \text{cyclic}),\qquad \omega_\cN=N_1N_2N_3,\qquad \omega_\cN S_i=-\half N_i,\qquad \textstyle\sum_iS_i^2=-\tfrac34\,\id :
\]
the $S_i$ are the rotations of this spatial triple and $\omega_\cN$ is its volume.

\begin{theorem}[One Pati--Salam family]
\label{thm:rung2}
\begin{enumerate}[label=(\roman*),leftmargin=2.2em,itemsep=1pt]
\item $\mathfrak g_2\cong\so6\cong\su4$. The generator of mass $G_2=\Gamma abcde$ is the volume element of $(\Gamma;P)$: a complex structure
that commutes with $\mathfrak g_2$ and with $\lint$ and separates the $\mathbf4$ from the $\bar{\mathbf4}$.
\item Of the sixteen left-handed Weyl fields of $\R^{64}$, eight are in the $\mathbf4$ and eight in the $\bar{\mathbf4}$.
\item The commutant of $\mathfrak g_2\oplus\lint$ has dimension $16$: it is $\C_\Omega\otimes\Cl(0,3)$, generated by $\Omega$ and
$N_1,N_2,N_3$. The volume $\omega_\cN$ is a symmetric involution of trace zero, and it splits $\Cl(0,3)=\mathbb H\oplus\mathbb H$. The
antisymmetric part of the commutant is $\uu2\oplus\uu2$: two commuting copies of $\su2$, spanned by $\half(\id+\omega_\cN)S_i$ and by
$\half(\id-\omega_\cN)S_i$, and the two phases $\half(\id\pm\omega_\cN)\Omega$.
\item Let $L$ be the complex space of left-handed Weyl multiplicities, of dimension sixteen. Under $\su4\oplus\su2\oplus\su2$ it carries
exactly $(\mathbf4,\mathbf2,\mathbf1)\oplus(\bar{\mathbf4},\mathbf1,\mathbf2)$. With the Cartan generators $H_1=\half ab$, $H_2=\half cd$
and $H_3=\half e\Gamma$ of $\mathfrak g_2$, $B\!-\!L=\pm\tfrac23(H_1+H_2+H_3)$ and $Q=T_{3L}+T_{3R}+\half(B\!-\!L)$, where $T_{3L}$ and
$T_{3R}$ are generators of the two $\su2$ of (iii): for one choice of the sign of $B\!-\!L$ and of the orientation that defines left-handed,
the sixteen left-handed fields have the charges $\tfrac23$ and $-\tfrac13$ three times each, $0$ and $-1$, and $-\tfrac23$ and $\tfrac13$ three
times each, $0$ and $1$.
\item $G_2=-\Omega\,\omega_\cN$ and $\Tr(\Omega G_2)=0$: the rotation generated by the generator of mass has one sense on eight left-handed
fields and the opposite sense on the other eight. The eight invariant kernels of \cref{thm:massrung}(v) commute with it, and under each of
the two $\su2$ they have spin one half.
\end{enumerate}
\checks{cf-2a to cf-2o, er-2k to er-2q; the dimensions in (iii) and the charges in (iv) numerically}
\end{theorem}

\begin{proof}[Proof of (iv), the family]
Complexify. The eigenspace $\Omega=+i$ of $\C^{64}$ is $S_L\otimes L$, the left-handed Weyl representation $S_L$ of $\lint$ tensored with
$L$, on which the gauge algebra acts. The six operators $(\Gamma;P)$ commute with $\lint$ and with $\Omega$, so they act on $L$, where they
generate the complex Clifford algebra of six generators, $M_8(\C)$; hence $L$ is two copies of its simple module $\C^8$. Under
$\mathrm{Spin}(6)\cong SU(4)$ each copy restricts to $\mathbf4\oplus\bar{\mathbf4}$, the eigenspaces of the volume element $G_2$ for the
eigenvalues $-i$ and $+i$, which fixes which one is called $\mathbf4$. (The real algebra generated by the six operators is
$\Cl(6,0)\cong M_4(\mathbb H)$ in our convention; only its complexification is used here.) The operator $\omega_\cN$ commutes with $\Omega$,
with $\lint$ and with $\mathfrak g_2$, and it is traceless on $L$, since $\omega_\cN$ and $\omega_\cN\Omega$ are products of distinct
generators. By (v), $G_2=-\Omega\,\omega_\cN$ acts on $L$ as $-i\,\omega_\cN$, so the $\mathbf4$ in $L$ is the eigenspace $\omega_\cN=+1$ and the
$\bar{\mathbf4}$ the eigenspace $\omega_\cN=-1$, each of dimension eight and each carrying two copies of its representation of $\su4$. On the
eigenspace $\omega_\cN=\pm1$ the generators $\half(\id\mp\omega_\cN)S_i$ of one $\su2$ vanish, and those of the other equal $S_i$, whose
Casimir $-\tfrac34$ is that of spin one half; this $\su2$ commutes with $\su4$ and acts on the multiplicity two as a doublet. Each eigenspace
is therefore one copy of $(\mathbf4,\mathbf2,\mathbf1)$, respectively $(\bar{\mathbf4},\mathbf1,\mathbf2)$. This proves the multiplicities,
the correlation of the $\mathbf4$ with one weak factor and of the $\bar{\mathbf4}$ with the other, and the absence of a mirror family. The
charges follow from the weights $(\pm\half,\pm\half,\pm\half)$ of the two spinors of $\so6$ under $(H_1,H_2,H_3)$, with an even and an odd
number of minus signs.
\end{proof}

This is the matter content of one family of the Pati--Salam model in its modern left--right form, with lepton number as the fourth
colour \cite{PatiSalam}, a right-handed neutrino included, with the right-handed fields
written as left-handed conjugates in the $(\bar{\mathbf4},\mathbf1,\mathbf2)$. The multiplet is established physics; what is specific here is
its realisation on the whole spinor module of the second rung under (A1), (A2) and (A4), with the spectator Lorentz action. Part~I obtains
the same content at the third rung with a chosen projector of rank $64$; here no projector is chosen. All three factors are gauged:
$\su4$ on the temporal directions by (A2), and $\su2\oplus\su2=\mathfrak s_2$ on the spatial triple of the second constituent by (A4), as its
two halves. The half that acts on the left-handed $\mathbf4$ is $\su2_L=\half(\id+\omega_\cN)S$. The charge of Part~I, $\half J=H_1$, is a
generator of $\mathfrak g_2$ and takes the values $\pm\half$ on every field.

By (v) the two senses of the rotation generated by $G_2$ distinguish the $\mathbf4$ from the $\bar{\mathbf4}$. In \PI{Section~7.1.2} the two
senses of the Dirac rest rotation are particle and antiparticle; here they distinguish the two halves of the family in the left-handed
convention, which by themselves do not classify particles and antiparticles. The eight invariant kernels are $\gamma^0$ times elements of
$\C_\Omega\otimes\Cl(0,3)$, and their count and spins are those of a bidoublet $(\mathbf1,\mathbf2,\mathbf2)$ of mass terms that pair the two
halves and do not break $\su4$.
\checks{cf-2h, cf-2m, for the count and the spins}

\begin{theorem}[The mass field at the second rung]
\label{thm:massfield2}
\begin{enumerate}[label=(\roman*),leftmargin=2.2em,itemsep=1pt]
\item The space $\mathcal M_2$ of Lorentz-invariant mass kernels has dimension $272=2\cdot136$, the real dimension of the complex symmetric
mass matrices of sixteen Weyl fields, $\mathrm{Sym}^2L$. Each of its elements is $C\omega_S$ with $C$ in the commutant $\Cl(6,3)$.
\item Under $\su4\oplus\su2\oplus\su2$,
\[
\mathcal M_2=2\,(\mathbf6,\mathbf1,\mathbf1)\oplus(\mathbf{15},\mathbf2,\mathbf2)\oplus(\mathbf{10},\mathbf3,\mathbf1)\oplus
(\overline{\mathbf{10}},\mathbf1,\mathbf3)\oplus(\mathbf1,\mathbf2,\mathbf2),
\]
of real dimensions $24$, $120$, $60$, $60$ and $8$. The labels are those of the fermion bilinears that the kernels multiply: with $\Omega$ as
the imaginary unit, a kernel $M$ carries under $M\mapsto[t,M]$ the charges of its bilinear, and the coefficient of the bilinear in the usual
complex notation transforms in the conjugate representation.
\item The Higgs-valued kernels $\sigma_2(V^+_2)\,\omega_S$ form a six-dimensional real subspace of the $24$-dimensional sextet part. The
Yang--Mills-valued kernels $\mathfrak g_2\,\omega_S$ form a fifteen-dimensional real subspace of the $120$-dimensional
$(\mathbf{15},\mathbf2,\mathbf2)$, and they are not invariant under the two $\su2$.
\item Projected on the four-plane, the two six-dimensional masses of \PI{Thm~9.62}(ii) land in different places: $T_2$ in the sextet part, a
vector of the gauged $\su4$, and $S_2$ in the $(\mathbf1,\mathbf2,\mathbf2)$, a singlet of $\su4$ that transforms under the two $\su2$. The
generator of mass times $\omega_S$, $G_2\,\omega_S=\pm\gamma^0\omega_\cN$, lies in the $(\mathbf1,\mathbf2,\mathbf2)$ as well.
\end{enumerate}
\checks{mr-1, mr-2, mr-3, mr-4; (ii) and (iii) numerically, with Casimir eigenvalues rounded to quarters}
\end{theorem}

\begin{proof}[Proof of (i) and (ii)]
A Lorentz scalar made from two left-handed Weyl fields is antisymmetric in the two spinor indices and hence symmetric in the internal ones,
so $\mathcal M_2$ is the real space underlying $\mathrm{Sym}^2L$, of real dimension $2\cdot\binom{17}2=272$. With
$L=(\mathbf4,\mathbf2,\mathbf1)\oplus(\bar{\mathbf4},\mathbf1,\mathbf2)$,
\[
\mathrm{Sym}^2L=(\mathbf{10},\mathbf3,\mathbf1)\oplus(\mathbf6,\mathbf1,\mathbf1)\oplus(\overline{\mathbf{10}},\mathbf1,\mathbf3)\oplus
(\mathbf6,\mathbf1,\mathbf1)\oplus(\mathbf{15},\mathbf2,\mathbf2)\oplus(\mathbf1,\mathbf2,\mathbf2),
\]
from $\mathrm{Sym}^2(V\otimes W)=\mathrm{Sym}^2V\otimes\mathrm{Sym}^2W\oplus\Lambda^2V\otimes\Lambda^2W$, $\Lambda^2\mathbf4=\Lambda^2\bar{\mathbf4}=\mathbf6$
and $\mathbf4\otimes\bar{\mathbf4}=\mathbf{15}\oplus\mathbf1$. The complex dimensions $30+6+30+6+60+4=136$ give the real dimensions stated.
In particular no part is a gauge singlet, and the only part on which $\su4$ acts trivially is the $(\mathbf1,\mathbf2,\mathbf2)$.
\end{proof}

These are the representations of the scalar fields of left--right symmetric Pati--Salam theories, in the decompositions of
\cite[eqs.~(3.5)--(3.7)]{DingKingLu2021}, which build on the construction of \cite{PatiSalam}: the
$(\mathbf1,\mathbf2,\mathbf2)$ for the Dirac masses of quarks and leptons, the right-handed triplet for the breaking to colour and
hypercharge with Majorana masses for the right-handed neutrinos, and the $(\mathbf{15},\mathbf2,\mathbf2)$ for splitting quark from lepton
masses. Under $\su4$ alone, all of $\mathcal M_2$ except the $(\mathbf1,\mathbf2,\mathbf2)$ transforms nontrivially; the
$(\mathbf1,\mathbf2,\mathbf2)$, which carries the projection of the spatial volume $S_2$, transforms only under the two $\su2$ of
$\mathfrak s_2$. With $\mathfrak s_2$ gauged, every part of $\mathcal M_2$ transforms nontrivially under the gauge algebra, as (A4) requires of
the mass.

\begin{theorem}[A Pati--Salam gauge theory]
\label{thm:ps}
\begin{enumerate}[label=(\roman*),leftmargin=2.2em,itemsep=1pt]
\item The antisymmetric part $\mathfrak s_2$ of the Clifford algebra of $N_1,N_2,N_3$ is spanned by the $\half N_i$ and the $\half N_jN_k$.
It is $\su2\oplus\su2$, the two halves $\half(\id\pm\omega_\cN)$ of the rotations of that triple. It commutes with $\mathfrak g_2$ and $\lint$ and keeps the kinetic
forms. The gauge algebra $\mathfrak g_2\oplus\mathfrak s_2=\su4\oplus\su2\oplus\su2$ has dimension $21$.
\item On the left-handed fields $\Tr(Y^3)=0$ for every $Y$ in the gauge algebra, and each $\su2$ has four doublets.
\item No nonzero element of $\mathcal M_2$ is invariant under the gauge algebra: there is no gauge-invariant mass term, and every fermion mass
needs a value of the mass field.
\item Colour is $\su3$, the derived algebra of the centraliser of $B\!-\!L$ in $\su4$. No nonzero element of the sextet part is invariant
under it: every value of the image of the clock vector breaks colour. The clock direction keeps $\so5$.
\item The mass field contains the directions of the Pati--Salam chain. In its $(\overline{\mathbf{10}},\mathbf1,\mathbf3)$ there is one
complex direction invariant under colour, under $\su2_L$ and under $Y=T_{3R}+\half(B\!-\!L)$, with stabiliser $\su3\oplus\su2\oplus\uu1$ of
dimension $12$. The generator of mass times $\omega_S$, $G_2\,\omega_S=\pm\gamma^0\omega_\cN$, lies in the $(\mathbf1,\mathbf2,\mathbf2)$, has
square $-\id$, and has stabiliser $\su4\oplus\su2_V$, of dimension $18$. Together they keep $\su3\oplus\uu1$, of dimension $9$, which contains
the electric charge $Q$. This holds for either labelling of the two $\su2$ as $L$ and $R$; which one is $L$ is fixed by the orientation that
defines the left-handed fields.
\end{enumerate}
\checks{ps-1, exactly; ps-2, ps-8, (ii) on random elements and exactly on the weights; ps-3 to ps-6, numerically with dimensions read as
ranks}
\end{theorem}

\begin{proof}[Proof of (ii)]
For $\su4$ the cubic anomaly of $L$ is $2A(\mathbf4)+2A(\bar{\mathbf4})=0$. A doublet has no cubic invariant, so each $\su2$ contributes
nothing, and the mixed traces between different simple factors vanish because each factor is traceless on its own representation. Hence
$\Tr Y^3=0$ for every $Y$; equivalently, every homogeneous cubic polynomial in the five Cartan coordinates $(H_1,H_2,H_3,T_{3L},T_{3R})$ sums to zero
over the sixteen weights, and conjugacy into a Cartan subalgebra extends this to the whole algebra. Each $\su2$ acts on four doublets, the
colour--lepton index of one half of $L$.
\end{proof}

The second rung is therefore a Pati--Salam gauge theory with one family \cite{PatiSalam}. Its perturbative gauge anomalies cancel, and each
weak factor has four Weyl doublets, an even number, so the theory passes the test of the $SU(2)$ global anomaly \cite{Witten82} for the
direct-product gauge group on spin four-manifolds. Other global forms are possible: the central element $(-\id_4,-\id_2,-\id_2)$ acts
trivially on the family, and a statement for the quotient group would need the bundles of that group. All fermion masses come from the mass
field of (A4). The two weak factors distinguish the $\mathbf4$ and $\bar{\mathbf4}$ sectors; their correlation makes the full Pati--Salam
representation chiral and forbids a gauge-invariant fermion mass. Neither algebra alone does this. The transported algebra alone is not chiral: the doublet of $\mathfrak g_1$ is
pseudo-real, $\mathfrak g_2$ acts on $\mathbf4\oplus\bar{\mathbf4}$ with equal multiplicities, and the spinor of $\mathfrak g_3$ is real. For
colour this is as it should be, the strong interaction treating left and right alike. In the Standard Model the chirality is carried by the
weak interaction, and here it is carried by the correlation of $\mathfrak s_2$ with $\su4$: on the left-handed fields the $\mathbf4$ are
doublets of one $\su2$ and the $\bar{\mathbf4}$ of the other. Nor is $\mathfrak s_2$ alone chiral: restricted to it the left-handed fields are
$4(\mathbf2,\mathbf1)\oplus4(\mathbf1,\mathbf2)$, and on either block the antisymmetric invariant of the doublet times an antisymmetric form
on the four copies is a symmetric invariant mass matrix. The two blocks admit $2\binom42=12$ complex invariant bilinears; imposing $\su4$
as well removes all of them, by (iii).
\checks{ps-10}

Two consequences concern the vacuum, and they describe directions; a potential or a vacuum assumption is needed to select them. First,
colour: the image of the clock vector lies in the sextet part, and any nonzero value of it breaks colour by (iv). No other vacuum value can
restore it, since for values in different isotypic parts the unbroken algebra is the intersection of their stabilisers. This is why, by
(A3), the image of the clock vector has no vacuum value at the second rung. Second, the generator of mass: carried into $V^-$ by $\omega_S$,
it gives the direction of the $(\mathbf1,\mathbf2,\mathbf2)$ that, together with a $(\overline{\mathbf{10}},\mathbf1,\mathbf3)$ direction,
breaks the gauge algebra to colour and electromagnetism by (v). The electroweak stage that follows is the subject of \cref{ch:ew}. With
$G_2\,\omega_S$ alone, whose square is $-\id$, the sixteen Weyl fields pair into eight Dirac fermions of equal mass, counting colours; the
masses are split by further directions of $\mathcal M_2$ (\cref{sec:split}).

\begin{proposition}[The vacuum family of Part~I inside $\mathfrak g_2$]
\label{prop:vacua2}
\begin{enumerate}[label=(\roman*),leftmargin=2.2em,itemsep=1pt]
\item $I_5X=G_2\,X\Gamma$ for every $X\in P$. So the family $\Pi(n)=\half(\id+I_5n)$ of \PI{Thm~8.58} is the spectral projector of the coset
generator $n\Gamma$ of $\mathfrak g_2$, with the generator of mass $G_2$ as the imaginary unit; at $n=e$ it is the vacuum $\half(\id+abcd)$
of \PI{Def.~8.56}.
\item The stabiliser of $\half(\id+abcd)$ in $\mathfrak g_2$ has dimension seven: the rotations of $(a,b,c,d)$ and $e\Gamma$,
$\so4\oplus\so2$. Restricted to $\mathrm{span}\{a,b,e\}$ the family is a vector of the $\so3$ on $(a,b,e)$, and its tangents at the vacuum
are the charged scalar of \PI{Thm~8.58}.
\item Inside $\mathfrak g_2$ the clock $\kappa\Gamma$ keeps $\so5$. The mass form $T$ of the first constituent and the projector $\half(\id+abcd)$ each keep
$\so4\oplus\so2$; at rung $n$ the mass form $T$ keeps an algebra of dimension $1+(3n-2)(3n-3)/2$.
\item $J_3=ab+cd+e\Gamma$ has the eigenvalues $\pm i$ and $\pm3i$ on the spinors, and $\half\,\mathrm{ad}(J_3)$ is a complex structure on the
six directions $(\Gamma;P)$. The centraliser of $J_3$ in $\mathfrak g_2$ has dimension nine, $\uu3=\su3\oplus\uu1$, and it contains the
charge $H_1=\half ab$ of Part~I. The electric charge $Q=T_{3L}+T_{3R}+\half(B\!-\!L)$ has components outside $\su4$ and is not an element of
this $\uu3$.
\end{enumerate}
\checks{cf-2n, cf-2g, sr-5, sr-6, sr-7, sr-8, sr-10, er-ns, in both forms}
\end{proposition}

None of the three vacua of (iii) keeps $\su3$: $\so4\oplus\so2$ has dimension seven, and $\so5\cong\spp2$ has no subalgebra $\su3$. A vacuum
along $J_3$ would keep $\su3\oplus\uu1$ inside $\mathfrak g_2$, the colour and $B\!-\!L$ of a Pati--Salam theory. These stabilisers are taken
inside $\mathfrak g_2$. In the full gauge algebra $\mathfrak g_2\oplus\mathfrak s_2$ the answer depends on the object held fixed, since passing
from an endomorphism to its mass kernel by $\omega_S$ changes which weak generators survive: the clock $\Gamma_2$ keeps an algebra of
dimension $13$ and its kernel $\Gamma_2\omega_S$ one of dimension $16$; the endomorphism $J_3$ keeps $15$ dimensions and the kernel $J_3\omega_S$
keeps $\uu3\oplus\su2_V$, of dimension $12$, whose weak factor is the diagonal one. These are assigned backgrounds; selecting one requires
dynamics or a vacuum assumption.
\checks{ps-11} Other triples of the form \cref{eq:triple} at the
second rung behave as at the first: with $X=E_7$ or the volume element $\Gamma_1$ of the first constituent the algebra is compact, commutes
with $\lint$ and keeps the kinetic forms, and its Casimir is $-\tfrac34$ on all of $\R^{64}$; for $X=E_7$ the module $\R^{64}$ is eight Weyl
doublets, an even number. The constituent's own clock $E_1$ fails, and a spatial direction gives no compact algebra.
\checks{sr-1, sr-2, sr-3, sr-4, sr-9}

\section{The third rung}

\begin{theorem}
\label{thm:rung3}
\begin{enumerate}[label=(\roman*),leftmargin=2.2em,itemsep=1pt]
\item $\mathfrak g_3\cong\so9$, of dimension $36$, and $\R^{512}$ carries $128$ left-handed Weyl fields. Six commuting complex structures,
four Cartan generators of $\so9$, $\Omega$ and one rotation of $\lint$, have traceless products, so each joint eigenspace has dimension
eight: the $128$ fields are eight copies of the spinor $\mathbf{16}$ of $\so9$. The commutant of $\mathfrak g_3\oplus\lint$ contains
$\C_\Omega\otimes\Cl(0,6)$, of dimension $128=\dim M_8(\C)$.
\item The directions $E_7,E_{13},E_{14},E_{15}$, the weak directions of Part~I, belong to $P$: their six bivectors lie in $\mathfrak g_3$.
\item The fifteen bivectors of $\cN$, the spatial directions of the second and third constituents, which contain the colour algebra of
Part~I, lie in the commutant.
\item The generator of mass $G_3$ is an involution, and $72$ of the blade kernels are invariant under $\mathfrak g_3$
(\cref{thm:massrung}); with $\mathfrak s_3$ as well a space of dimension two remains (\cref{prop:rung3s}).
\end{enumerate}
\checks{er-3k, er-3l, er-3m, er-3n, cf-3e, cf-3f, cf-3h}
\end{theorem}

The $\mathbf{16}$ of $\so9$ is the restriction of the spinor of $\so{10}$, which carries one family with a singlet neutrino
\cite{FritzschMinkowski}, \cite[\S93.2.3]{PDGGUT}; the restriction through $\so9$ is our calculation. Under
$\so6\oplus\so3\subset\so9$ it is $(\mathbf4,\mathbf2)\oplus(\bar{\mathbf4},\mathbf2)$.

\begin{proposition}[The third rung under (A4)]
\label{prop:rung3s}
At rung three (A4) gauges $\mathfrak s_3$, the antisymmetric part of the Clifford algebra of the six spatial generators of the second and
third constituents. It has dimension $28$, commutes with $\mathfrak g_3$ and $\lint$, contains the fifteen colour bivectors of Part~I, and its
adjoint representation is irreducible: $\mathfrak s_3=\so8$. Under $\mathfrak g_3\oplus\mathfrak s_3=\so9\oplus\so8$ the space of left-handed
Weyl multiplicities is the irreducible $(\mathbf{16},\mathbf8)$. The kernel $M=G_3\,\omega_S$ is a Lorentz-invariant mass kernel, antisymmetric
and invariant under $\mathfrak g_3\oplus\mathfrak s_3$, with $M^2=-\id$: for every $m\ne0$, $mM$ has full rank, so the gauge symmetry permits
a common nonzero mass for every mode, and the theory $\so9\oplus\so8$ at rung three is not chiral. The kernels invariant under $\mathfrak g_3\oplus\mathfrak s_3$
form the real span of $M$ and $M\Omega$: for every such kernel $K$, $M^{-1}K$ lies in the joint commutant $\C_\Omega$. Both are
antisymmetric, square to $-\id$ and anticommute, so $aM+bM\Omega$ with $(a,b)\ne(0,0)$ gives the same singular value $\sqrt{a^2+b^2}$ to
every mode. The two real coordinates count invariant matrices before phase identifications, and the common mass scale is a free
coefficient, since (A3) selects no vacuum at this rung.
\checks{ps-7, ps-9, fm-4}
\end{proposition}

The rule (A4), carried to the third rung, therefore gauges the colour of Part~I but gives a representation that is real and admits a
Majorana-type mass, unlike the pseudo-real single doublet of the first rung, which admits no invariant bilinear mass. Nothing gauged on
$\cN$ changes this, since $G_3\,\omega_S$ commutes with $\mathfrak g_3$ as well. Read as an interaction (\cref{sec:interactions}), this is a
property of the third rung: chirality is required of the weak interaction, which the second rung carries. What interaction the third rung
carries is open (\cref{sec:open}).

\begin{remark}[Counts]
The spinor of $\so{3n}$ has dimension $2^{\lfloor3n/2\rfloor}$; for even $n$ this is the full complex spinor, whose two half-spinors have
dimension $2^{3n/2-1}$ each. $\R^{8^n}$ carries $2^{3n-2}$ left-handed Weyl fields. For odd $n$ they
are $2^{(3n-3)/2}$ copies of the spinor: one doublet at $n=1$ and eight $\mathbf{16}$ at $n=3$. For even $n$ they are $2^{3n/2-2}$ copies of
each of the two chiral spinors: two $\mathbf4$ and two $\bar{\mathbf4}$ at $n=2$.
\end{remark}

\section{The third rung as a family index}
\label{sec:families}

Climbing from the second rung to the third adds a constituent. Its spinor index can be read as a family index, and what that index carries
depends on which algebra is taken to act.

\begin{theorem}[Families at the third rung]
\label{thm:families}
Write $\R^{512}=\R^{64}\otimes\R^8$, the second factor the spinors of the third constituent, with its pole operators
$K_j=E_{12+j}E_{15+j}$, its Fock number $N=\half(3\,\id-\sum_jK_j)$ and its diagonal rotations $\half(e_ie_j-e_{i+3}e_{j+3})$, which turn its
two triples together.
\begin{enumerate}[label=(\roman*),leftmargin=2.2em,itemsep=1pt]
\item With the interaction of the second rung held on the first two constituents, its algebras built with $\Gamma_2$, the $128$
left-handed Weyl fields are eight copies of the family of \cref{thm:rung2}, and there is no mirror family
$(\mathbf4,\mathbf1,\mathbf2)\oplus(\bar{\mathbf4},\mathbf2,\mathbf1)$.
\item $N$ and the diagonal rotations commute with that interaction, with $\lint$ and with $\Omega$. Under them the eight copies are
$\mathbf1\oplus\mathbf3\oplus\mathbf3\oplus\mathbf1$ at $N=0,1,2,3$: the triplets are the octants of weight one and two, where Part~I places
$\bE$ and $\bB$ (\PI{Introduction and Section~3.5}).
\item Restricted to the embedded Pati--Salam algebra built with $\Gamma_3$, the rotations of $(\Gamma_3;a,b,c,d,e)$ together with
$\mathfrak s_2$, the $128$ fields are four families and four mirror families: the families at $N=1$ and $3$, the mirrors at $N=0$ and $2$.
The Fock parity is $(-1)^N=K_1K_2K_3$, which is minus the ordered volume of the third constituent.
\item The full third-rung algebra $\mathfrak g_3\oplus\mathfrak s_3$ does not preserve the Fock grading: the generator $Y=\half E_2E_{13}$ of
$\mathfrak g_3$ anticommutes with $K_1$, commutes with $K_2$ and $K_3$, and satisfies $[N,Y]=\half E_2E_{16}\ne0$; it changes the Fock parity.
Under $\so9\oplus\so8$ the Weyl multiplicities form the irreducible $(\mathbf{16},\mathbf8)$ of \cref{prop:rung3s}, in which the families
and mirrors of (iii) are mixed.
\end{enumerate}
\checks{fm-1, fm-2, fm-3, fm-4, numerically; the identities of (iii) and (iv) exactly}
\end{theorem}

\begin{figure}[t]
\centering
\begin{tikzpicture}[font=\footnotesize,>=Stealth,
  fam/.style={circle,draw,fill=blue!25,minimum size=8mm,inner sep=0pt},
  mir/.style={circle,draw,fill=white,minimum size=8mm,inner sep=0pt}]
\node[mir] (o) at (0,0) {$000$};
\node[fam] (a1) at (-1.8,1.4) {$100$}; \node[fam] (a2) at (0,1.4) {$010$}; \node[fam] (a3) at (1.8,1.4) {$001$};
\node[mir] (b1) at (-1.8,2.8) {$110$}; \node[mir] (b2) at (0,2.8) {$101$}; \node[mir] (b3) at (1.8,2.8) {$011$};
\node[fam] (f) at (0,4.2) {$111$};
\foreach \u/\v in {o/a1,o/a2,o/a3,a1/b1,a1/b2,a2/b1,a2/b3,a3/b2,a3/b3,b1/f,b2/f,b3/f} \draw[black!45] (\u) -- (\v);
\node[anchor=west] at (3.0,0) {$N=0$: spin $0$, the vacuum octant};
\node[anchor=west] at (3.0,1.4) {$N=1$: spin $1$, the octants of $\bE$};
\node[anchor=west] at (3.0,2.8) {$N=2$: spin $1$, the octants of $\bB$};
\node[anchor=west] at (3.0,4.2) {$N=3$: spin $0$, the full octant};
\node[fam,minimum size=4mm] (k1) at (-5.6,3.2) {}; \node[anchor=west] at (-5.3,3.2) {family};
\node[mir,minimum size=4mm] (k2) at (-5.6,2.6) {}; \node[anchor=west] at (-5.3,2.6) {mirror family};
\node[anchor=west,align=left,text width=3.6cm] at (-5.8,1.2) {colours: the restriction (iii);\\ in reading (i) all eight\\ octants carry families};
\end{tikzpicture}
\caption{The eight octants of the third constituent, labelled by their excited Fock bits and arranged by the Fock number $N$
(\cref{thm:families}). With the interaction of the second rung held, every octant carries one family, and the diagonal rotations group them
as $\mathbf1\oplus\mathbf3\oplus\mathbf3\oplus\mathbf1$. Under the Pati--Salam algebra built with $\Gamma_3$ the octants of odd weight carry
families and those of even weight mirror families; the full third-rung algebra $\so9\oplus\so8$ mixes the two kinds.}
\label{fig:octants}
\end{figure}

Climbing therefore gives a family index, but counting does not give three. With the interaction of the second rung held there are eight
chiral copies, and at higher rungs $8^{n-2}$; on restriction to the embedded Pati--Salam algebra built with $\Gamma_3$ there are four
families and four mirrors, which the full third-rung algebra mixes. The three that the framework offers is a triplet of the Fock grading,
the three directions of a triple (\cref{fig:octants}). The weight-one triplet is a set of families in readings (i) and (iii), the weight-two
triplet only in (i). The routes to three families are collected, and left open, in \cref{sec:open}.
